\pagenumbering{roman}   % front matter; Chapter 1 starts at arabic page 1
\begin{titlepage}
\thispagestyle{empty}
\vspace*{3.5cm}
\noindent\rule{\textwidth}{2pt}
\begin{center}
  {\LARGE\bfseries Offline RL: From Theory to Industrial Practice\par}
  \vspace{0.6em}
  {\large Practical notes on offline deep reinforcement learning\par}
\end{center}
\noindent\rule{\textwidth}{0.6pt}

\vspace{1.6cm}
\begin{center}
  {\large\bfseries Artem Zvyagintsev\par}
  \vspace{0.3em}
  {\normalsize\url{https://corba777.github.io/offline-rl-book/}\par}
  \vspace{1.2em}
  {\normalsize\scshape A Preprint\quad--\quad Version $version$\quad\textperiodcentered\quad $date$\par}
\end{center}

\vspace{1.6cm}
\begin{center}\textbf{Abstract}\end{center}
\begin{quote}
Offline reinforcement learning trains a control policy from data that has
already been logged, with no experiments on the real system. It matters wherever
trial and error is expensive or unsafe, as on an industrial plant.

The book explains why the problem is hard and how it is solved. It starts with
simply imitating the operator and with the central trap: the learned values
overrate actions the data never shows. It then covers evaluating a policy from
logs and the main families of methods: conservative value estimates, policy
constraints, Decision Transformers and learning with a model of the
environment. The last chapters turn to industrial practice: building process
physics into the model, a larger case study checked over many runs, explaining
decisions, and LLM agents.

Each chapter gives the idea, the mathematics, working code and an honest look
at where the method breaks down. The code is for learning: it shows how the
algorithms work and is not meant as a production library.
\end{quote}

\vfill
\begin{center}\footnotesize
DOI: \href{https://doi.org/10.5281/zenodo.24045871}{10.5281/zenodo.24045871}\quad Text and figures: CC BY 4.0.\quad Code: MIT.\\
Source: \url{https://github.com/corba777/offline-rl-book}\\
Built from commit $commit$.
\end{center}
\end{titlepage}
\setcounter{page}{2}   % the title page is i, the contents start at ii
